\begin{theorem}[A finite cycle criterion when the cohomology is even]
\label{c:finite-even-criterion}
\status{Proved in the pure-even case}
Let $X/\C$ be smooth projective and connected, with $H^{\rm odd}(X,\Q)=0$,
let $n=\dim X>0$, and let $o\in\CH^n(X)_\Q$ have degree $1$.
Put $d=\dim_\Q H^*(X,\Q)=\chi(X)$.  There is a finite list of rational
tensor expressions $P_1,\ldots,P_s$, constructed by the procedure below
from the rational cup-product table and trace of $X$, for which the
following criterion is necessary and sufficient for all the relations
of Theorem~\ref{c:matrix-presentation} to hold in rational Chow groups
on every power:
\begin{enumerate}[label=(\roman*)]
\item $c_n(T_X)=d\,o$ and $\Delta_*o=o\boxtimes o$;
\item $\bigwedge^{d+1}h(X)=0$ in rational Chow motives;
\item $P_a(h(X),\mu,1,\operatorname{tr},o,\beta^{-1})=0$ for
      $1\leq a\leq s$.
\end{enumerate}
Here $\mu$ is the small-diagonal multiplication,
$\operatorname{tr}:h(X)\to\Q(-n)$ is the trace,
$\beta=\operatorname{tr}\circ\mu$, and
$\beta^{-1}:\Q(-n)\to h(X)\otimes h(X)$ is its inverse pairing.
Every condition is the vanishing of a specified finite linear
combination of correspondence diagrams.  Subject to (i), each can
equivalently be written as a finite list of cycles in the requested
point-and-diagonal rings.  The construction uses finite-dimensional
polynomial algebra, so is a finite criterion, rather than one separate
test for each power of $X$.
\end{theorem}

\begin{proof}
Write $A=H^*(X,\Q)$ as an ordinary vector space; the pure-even
hypothesis makes its tensor symmetry the ordinary one.  Retain the
Tate degree by a one-dimensional auxiliary representation $L$, and
put $G=\operatorname{GL}(A)\times\mathbb G_m$, with the second factor
acting on $L$.  The data
\[
 t=(\mu,1,\operatorname{tr},o,\beta^{-1})
\]
form one rational point of the finite-dimensional rational
$G$-representation
\[
 W=(A\otimes A^\vee\otimes A^\vee)\oplus A
   \oplus(A^\vee\otimes L^n)
   \oplus(A\otimes L^{-n})
   \oplus(A\otimes A\otimes L^{-n}).
\]
In the point $t$, the coordinate denoted by $o$ is the cohomology class
of the degree-one zero-cycle.  Let $I\subset\Q[W]$ be the ideal of
the Zariski closure of $G\cdot t$.  We describe all steps producing
the finite list.

First compute generators of $I$ by elimination: write the orbit map
in matrix entries of $g\in\operatorname{GL}(A)$ and a nonzero scalar
$z$, introduce inverses with equations $u\det(g)=1$ and $vz=1$, and
eliminate $g,u,z,v$ from the graph ideal.  The resulting elimination
ideal is the ideal of the image closure.  Hilbert's basis theorem
makes the resulting list finite.  Enlarge its span to a finite-dimensional
$G$-stable subspace $E\subset I$ generating $I$ as an ideal.
This enlargement is finite because the action on the polynomial ring
is locally finite and preserves the finite degree bound of the chosen
generators.  It can be computed by closing their span under the matrix
units of the Lie algebra and the second torus action; $G$ is connected
in characteristic zero.

The inclusion $E\hookrightarrow\Q[W]$ is equivalently a polynomial
covariant
\[
 p_E\in(E^\vee\otimes\Q[W])^G.
\]
Choose a split $G$-embedding of $E^\vee$ into a finite direct sum $T$
of mixed tensor powers of $A,A^\vee,L,L^\vee$.  Such a split embedding
exists because rational representations of $G$ are semisimple and
are generated by its defining representations and their duals.
For an explicit terminating search, enumerate these finite direct
sums and compute the rational intertwiner spaces.  Enumerate their
rational elements until one has full rank on $E^\vee$, and then
solve the linear intertwiner equations for a retraction.  The
existence of a split embedding proves termination.  The components of the image of
$p_E$ in $T\otimes\Q[W]$ are the covariants we denote by $P_a$.

We recall the two representation-theoretic facts needed here.
Over $\C$, averaging a Hermitian form over
$U(d)\times U(1)$ gives invariant complements.  This compact group
is Zariski dense in $G_\C$, so the complements are $G_\C$-stable.
The intertwiner and splitting equations have rational coefficients;
Gaussian elimination descends their solvability to $\Q$.
Every finite-dimensional rational representation embeds, by its
matrix coefficients, into finitely many copies of $\Q[G]$.
The relevant coefficient functions have bounded polynomial degree
and bounded powers of $\det^{-1}$ and the second torus coordinate.
They therefore occur in mixed tensor powers, using symmetric and
exterior projectors.  Complete reducibility makes the embedding
split.  This proves the existence statements just used.  It also
proves exactness of invariants for the locally finite rational
modules below: place a preimage and its translates in a
finite-dimensional submodule and split the resulting surjection.
The exact references are \citet[Theorem~4.14, pp.~88--89;
Theorem~22.42, pp.~478--479; Remark~12.53(a), p.~248,
and the Reynolds operator on p.~249]{c_milneag2022}.

We explain why these covariants are tensor expressions of the promised
kind.  Polarize a homogeneous polynomial covariant in the finitely
many tensor coordinates of $W$.  Its coefficients are invariants in
mixed tensor powers of $A$ and $A^\vee$.  The first fundamental
theorem says that these are sums of complete contractions, with any
uncontracted indices supplying the target tensor.  The only linear
relations among such contractions are generated by the
$(d+1)$-fold antisymmetrizer and by the scalar loop value $d$.
These are precisely \citet[Theorems~1.2--1.5, pp.~4--5]{c_procesi2021tensor};
polarization and restitution apply because the characteristic is zero.
This use is of the classical tensor theorems, and does not apply a
matrix-only statement to arbitrary tensor variables without that step.
The extra $\mathbb G_m$ factor merely selects the matching Tate
degrees.  Thus each $P_a$ is a finite rational linear combination of
the prescribed diagrams, and every universal identity in their
ordinary $d$-dimensional interpretation follows from the stated
antisymmetrizer and loop relation.

Now let $F$ be any polynomial covariant with target a mixed tensor
space $T'$ that vanishes at $t$.  Equivariance implies that it
vanishes on $G\cdot t$ and its closure.  Hence
\[
 F\in(T'\otimes I)^G.
\]
The multiplication map $E\otimes\Q[W]\twoheadrightarrow I$ is
$G$-equivariant.  Since $G$ is reductive over $\Q$, taking invariants
is exact, and $F$ lifts to
\[
 \widetilde F\in(T'\otimes E\otimes\Q[W])^G.
\]
Contracting its $E$ index against $p_E$ gives $F$ exactly.  Extend
that contraction through the chosen split embedding into $T$.
The first fundamental theorem again writes $\widetilde F$ as a
diagram.  Consequently the tensor ideal of all relations at $t$
is generated by the finite covariants $P_a$, together with the
$(d+1)$-fold antisymmetrizer and the loop relation.  No reductivity
of the stabilizer of $t$ has been assumed or used.

Evaluate this presentation in Chow motives.  The scalar loop value
is $\chi(X)=d$ by the categorical trace formula.  Condition (ii)
annihilates the dimension relation, and condition (iii) annihilates
the finite orbit-ideal generators.  Therefore every cohomological
diagram relation vanishes in Chow, in particular every relation in
the requested rings.  This proves sufficiency.
The additional identities for the auxiliary line,
$\bigwedge^2 L=0$ and $\operatorname{tr}(\id_L)=1$, already hold
for the fixed even invertible Tate object $L=\Q(-1)$; they are
not additional conditions on $(X,o)$.

Conversely suppose all the requested relations hold in Chow.
The cohomological relations
$\Delta\cdot o_1=o_1o_2$ and
$\Delta^2=d\,o_1o_2$ hold there.  The first is the second assertion
of (i); pushing the second to $X$ gives its first assertion.
The alternating projector in (ii) is a rational linear combination
of permutation graphs on $X^{d+1}$, each a product of pairwise
diagonals in $X^{2(d+1)}$.  Its cohomology is zero, so (ii) follows.

Finally the correspondence diagrams used above do not enlarge the
vanishing test after (i).  The Frobenius identities reduce each
connected diagram to a small diagonal with some point insertions
and some handles.  A handle inserts $c_n(T_X)$; a point insertion
inserts $o$.  By (i) both reduce to scalar multiples of $o$, two
such insertions vanish for dimension reasons, and a single insertion
on a small diagonal is the external product of the point classes.
Closed components give their degrees.  Thus every diagram reduces
to the pointed partitions of Definition~\ref{c:def-pointed}.
This is the same graph reduction as Theorem~\ref{c:graph-rule},
now including contractions and using the two verified Chow
identities.  Each $P_a$ has zero cohomology and reduces to a relation
of the requested kind; it consequently vanishes in Chow.  This
proves necessity and the finite criterion.
\end{proof}

\begin{remark}[Scope of the finite criterion]
\label{c:finite-criterion-scope}
\status{Proved scope clarification; the arbitrary-pair classification remains unresolved}
Theorem~\ref{c:finite-even-criterion} applies to arbitrary cup-product
algebras with pure-even cohomology, not only to surfaces.  It gives
a finite algebraic construction of the cycle conditions, without
claiming to decide their rational equivalence by computation.
For mixed cohomology the direct use of
$\operatorname{GL}_{a|b}$ would require exactness of its invariant
functor, which fails.  Theorem~\ref{c:finite-super-criterion}
replaces that step by projective lifting and an external algebraic
parity correspondence.  Corollary~\ref{c:finite-weil-criterion}
also separates finite existence over an arbitrary Weil coefficient
field from effective construction with specified coefficient data.
Theorem~\ref{c:finite-unconditional-criterion} subsequently removes
the sign hypothesis by first passing to scalar forms in the
covariant polynomial category.  These finite criteria do not
decide the resulting Chow vanishings.
The unrestricted classification of all $(X,o)$ therefore remains
distinct from the specific K3 obstruction proved below.
\end{remark}

